
\vspace{5mm}
\begin{ledgroup}
\footnotesize
\pstart
\noindent
Die vorliegenden zwei Stücke weisen eine gemeinsame Entstehungsgeschichte auf:
Das Konzept N.\;\ref{58293}, 
dessen Wasserzeichen für den Zeitraum Oktober, unter anderem im Konzept N.\;\ref{58300}, gut belegt ist,
bezeugt Leibnizens (in den 1690er Jahren häufige) Beschäftigung mit den Regeln des schiefen Mehrkörperstoßes und ihren Folgen für die Krafterhaltung.
Auf einem separaten Zettel hat Leibniz einige Hilfsrechnungen zu einem besonders schwierigen Teil von N.\;\ref{58293} festgehalten (N.\;\ref{39763}).
Ein weiterer auf diesem Träger fragmentarisch überlieferter Text bietet einen Terminus post quem 
für die (als etwa simultan anzunehmende) Entstehung unserer Stücke 
und ermöglicht ihre Datierung in der Zeitspanne Mai 1693 bis 1696.
\pend
\pstart
Die Zusammengehörigkeit beider Stücke wird durch einen Vergleich der (Hilfs-)Rechnungen in N.\;\ref{39763}  mit den Art. (36)\textendash(38) von N.\;\ref{58293} ersichtlich.
So entspricht die Passage auf \refpassage{04_01_04d_02_004v_2a}{04_01_04d_02_004v_2b} von N.\;\ref{39763} 
den gestrichenen Art.~(36)\textendash(37) von N.\;\ref{58293}:
\refpassage{35_13_03_220-221_20a}{35_13_03_220-221_20b}
(das Verhältnis \textit{r} wurde in Art.~(35) eingeführt:
\refpassage{35_13_03_220-221_19a}{35_13_03_220-221_19b})
und die Passage auf \refpassage{04_01_04d_02_004v_1a}{04_01_04d_02_004v_3} von N.\;\ref{39763}
dem Art.~(38) von N.\;\ref{58293} (\refpassage{35_13_03_220-221_21a}{35_13_03_220-221_21b})
sowie der gestrichenen Randanmerkung zu
\refnote{FnA:35_13_03_220-221_rechnung}.
\pend
\pstart
Die kurze Aufzeichnung N.\;\ref{39763} ist auf einem Zettel mit der Signatur LH~IV~1, 4d, 2~Bl.~4 überliefert.
Auf Bl.~4~v\textsuperscript{o} befindet sich ein durch Papierabriss und -beschnitt nur noch fragmentarisch erhaltener Briefanfang; 
Bl.~4~r\textsuperscript{o} überliefert eine (wohl vollständige) Aufzeichnung über \indexN{\textls{Descartes} (Cartesius, des Cartes), René 1596\textendash1650}Descartes (\textit{Mons.~des Cartes a fait comme les charlatans}...), die die Seite ausfüllt. 
(Sie liegt in \cite{02150}\textsc{Foucher de Careil}, \textit{Nouvelles Lettres et Opuscules Inédits de Leibniz}, Paris 1857, S.\;12 gedruckt vor und wird vsl. in einem Band der Sechsten Reihe erscheinen.)
Dieser Umstand legt die Vermutung nahe, dass Leibniz erst den Briefentwurf verfasst habe, dann die Rückseite des Briefs für die Descartes-Aufzeichnung wiederverwendet und anschließend den Zettel auf diese zugeschnitten;
erst an dritter Stelle dürfte die Aufzeichnung N.\;\ref{39763} entstanden sein, für die Leibniz offenbar den freigebliebenen Rand von Bl.~4~v\textsuperscript{o} und sogar Zwischenräume im beschriebenen Teil verwendet hat.
Der Briefentwurf bietet einen Terminus post quem für die Entstehung der Notiz über Descartes und diese, die allerdings noch nicht verlässlich datiert ist, wiederum für die Entstehung der Aufzeichung N.\;\ref{39763}.
Der (für unsere Aufzeichung N.\;\ref{39763} also datierungsrelevante) Briefentwurf diente als Konzept zu einem nicht erhaltenen Schreiben Leibnizens an
\indexN{\textls{Brosseau}, Christophe 1630\textendash1717}Christophe Brosseau; auf dieses nahm \indexN{\textls{Brosseau}, Christophe 1630\textendash1717}Brosseau in einem \cite{02148}Brief vom 18. Mai Bezug (\textit{LSB} I, 9 N.\;276; siehe die Erläuterung, S. 437\,f.).
Sowohl in dem vorliegenden Entwurf als auch in Brosseaus Antwort geht es um die Frage der korrekten Abrechnung gewisser Zeitschriften ("gazettes"), nämlich auf Leibnizens Kosten statt auf die des 
\indexN{\textls{Braunschweig-Lüneburg}, Ernst August von, Herzog und Kurfürst von Hannover 1679\textendash1698}Kurfürsten von Braunschweig-Lüneburg. 
Leibniz bezeichnet im Briefentwurf seinen Dienstherren zunächst mit der Abkürzung "S.\,A.\,S.", ändert diese jedoch nachträglich in "S.\,A.\,E." (\textit{Son Altesse Électorale}), gemäß der am 19. Dezember 1692 verliehenen Kurwürde.
Im verschollenen Brief vom 8. Mai, sowie in einem Schreiben an die Kammer (Entwurf: LBr 560 Bl.~35), bemühte sich Leibniz um eine Klärung des Vorfalls (siehe \textit{LSB} I,~9, S.\;401). 
Aufgrund dieser Identifizierung ist der Briefentwurf auf dem Träger der Aufzeichnung N.\;\ref{39763}, Bl.~4~v\textsuperscript{o}, auf den 8. Mai 1693 zu datieren; dadurch ist ein Terminus post quem für N.\;\ref{39763} selbst gegeben, der wiederum bei der Datierung des Konzepts N.\;\ref{58293} berücksichtigt wird. 
\pend
\end{ledgroup}
